\documentclass[11pt]{article}

\usepackage[margin=1in]{geometry}
\usepackage{amsmath}
\usepackage{booktabs}
\usepackage{siunitx}
\usepackage{graphicx}
\usepackage{float}
\usepackage{fancyhdr}
\usepackage[hidelinks]{hyperref}

\setlength{\parindent}{0pt}
\setlength{\parskip}{0.5em}

% ---------------------------------------------------------------- Page header
\pagestyle{fancy}
\fancyhf{}
\lhead{June 5, 2026}
\rhead{Gas Injury Detection --- Progress Report}
\cfoot{\thepage}
\renewcommand{\headrulewidth}{0.4pt}

\begin{document}

% ---------------------------------------------------------------- Title block
\thispagestyle{empty}
\begin{center}
    {\Huge\bfseries Gas Injury Detection}\\[0.5em]
    {\large Wearable Sensor Classification --- Progress Report}\\[4em]
    {\large June 5, 2026}
\end{center}

\vspace{2em}

% ---------------------------------------------------------------- Key Results
\section*{Key Results}

\begin{itemize}
    \item \textbf{The current sensor network reliably detects blood mixtures down to
          20\% blood.} 

    \item The best model, \textbf{XGBoost with feature reduction and causal temporal
          smoothing}, runs in real time and reaches \textbf{95.8\% accuracy} on held-out
          sessions, with the few remaining errors confined to the moment a sample is
          introduced.

    \item \textbf{NH\textsubscript{3} is the dominant blood marker} --- its level scales
          with blood concentration, consistent with ammonia release from protein
          breakdown in blood.

\end{itemize}

% ---------------------------------------------------------------- Detection floor
\section*{Detection Floor: Why \texttt{mix\_140\_10} Was Excluded}

\texttt{mix\_140\_10} is the lowest-concentration mixture: 6.7\% blood in 93.3\% sweat.
NH\textsubscript{3} is the primary discriminator between blood and sweat, and its response
scales with blood content. At 6.7\% the NH\textsubscript{3} response falls into the
pure-sweat range, and two of the five \texttt{mix\_140\_10} sessions showed no measurable
NH\textsubscript{3} elevation above sweat on either the peak or the sustained level.

Since the two classes overlap in the sensor reading, the model cannot separate them, and
including \texttt{mix\_140\_10} only adds noise to the blood class, so these sessions were
excluded. The next ratio up, \texttt{mix\_120\_30} at 20\% blood, produces roughly double
the sweat NH\textsubscript{3} and separates cleanly, which sets 20\% as the lowest blood
content the network can reliably detect.

\begin{figure}[H]
    \centering
    \includegraphics[width=0.7\textwidth]{mix140_overlap.png}
    \caption{Sustained NH\textsubscript{3} per session by sample type. The 6.7\% blood
    mixture (\texttt{mix\_140\_10}, green) interleaves with the pure-sweat range (shaded),
    while the 20\% mixture (\texttt{mix\_120\_30}, blue) sits clearly above it. This
    overlap is why 6.7\% samples cannot be added to the model.}
\end{figure}

% ---------------------------------------------------------------- Models
\section*{Models Evaluated}

Three model families were compared, chosen to span interpretability, temporal-modeling
capacity, and suitability for embedded deployment.

\begin{table}[H]
\centering
\caption{Model characteristics.}
\begin{tabular}{p{2.3cm}p{5.3cm}p{5.3cm}}
    \toprule
    \textbf{Model} & \textbf{Strengths} & \textbf{Limitations} \\
    \midrule
    \textbf{Random Forest} &
    Interpretable feature importances; robust on small tabular datasets; no input scaling
    required. &
    Requires hand-crafted window aggregates; does not learn from the raw time-series. \\
    \addlinespace
    \textbf{XGBoost} &
    Gradient boosting corrects residual errors sequentially, outperforming RF on the same
    features; L1/L2 regularization in the training objective penalizes model complexity
    directly. &
    Same tabular feature dependency as RF; less interpretable. \\
    \addlinespace
    \textbf{1D CNN} &
    Learns temporal patterns directly from the raw signal; no manual feature engineering. &
    Data-hungry; input normalization suppressed the low-magnitude NH\textsubscript{3}
    signal, so it confused sweat with blood. \\
    \bottomrule
\end{tabular}
\end{table}

The Random Forest and XGBoost operate on 60\,s window aggregates (mean, std, max, slope)
of the engineered feature set. The CNN consumes the raw 30\,s window directly. Across
held-out sessions the tree models were clearly stronger: the task is fundamentally one of
signal \emph{level} (how high NH\textsubscript{3} rises), which the raw feature values
expose directly, whereas the CNN's per-channel scaling compressed that level away.

% ---------------------------------------------------------------- Dataset
\section*{Dataset}

Data was collected across 28 sessions using the ESP32-based sensor array (VOC,
NH\textsubscript{3}, HCHO, H\textsubscript{2}S, EtOH, plus temperature and humidity). The
6.7\% mixture (\texttt{mix\_140\_10}) is excluded for the reason described above. Each
session is manually labeled into baseline, sweat, and blood regions.

Splits are performed at the session level so that all windows from a given session stay
together, which prevents temporally correlated windows from leaking between training and
test. 24 sessions are used for training and 4 are held out as unseen test data.
Non-overlapping 60\,s windows produce 1{,}683 windows in total; the training windows are
class-balanced (sweat bootstrapped, baseline undersampled) before the tree models are fit.

\begin{table}[H]
\centering
\caption{Dataset summary.}
\begin{tabular}{lr}
    \toprule
    \textbf{Property} & \textbf{Count} \\
    \midrule
    Total sessions & 28 \\
    \quad Sweat & 8 \\
    \quad Pure blood & 10 \\
    \quad Mixture, 50\% blood (\texttt{mix\_75\_75}) & 5 \\
    \quad Mixture, 20\% blood (\texttt{mix\_120\_30}) & 5 \\
    \addlinespace
    Training sessions & 24 \\
    Held-out test sessions & 4 \\
    \addlinespace
    Total windows (60\,s) & 1{,}683 \\
    Test windows & 215 \\
    \bottomrule
\end{tabular}
\end{table}

% ---------------------------------------------------------------- Balancing
\section*{Training Set Balancing}

The 1{,}468 raw training windows are heavily imbalanced. Baseline dominates because each
session spends most of its time before the sample is introduced, and sweat is the minority
class. Training on this distribution would bias the model toward baseline and
underrepresent sweat, so the three classes are balanced to the blood count of 442 windows.

\textbf{Bootstrapping the sweat class.} Each synthetic sweat window is created by drawing
an existing sweat window at random (with replacement) and adding Gaussian noise to every
feature, scaled to 5\% of that feature's standard deviation across the sweat windows:
$x_{\text{syn}} = x + \mathcal{N}(0, \sigma^2)$ with $\sigma = 0.05 \times \text{std}(x)$.
The noise keeps each synthetic window recognizably sweat while avoiding exact duplicates,
which would only reweight the same points without adding variety. This adds 282 synthetic
windows, taking sweat from 160 to 442.

\textbf{Undersampling the baseline class.} Baseline is randomly subsampled from 866 down
to 442 windows.

The final balanced training set is 1{,}326 windows, 442 per class. Balancing is applied to
the training data only; the held-out test sessions keep their natural class distribution.

\begin{table}[H]
\centering
\caption{Training window counts before and after balancing.}
\begin{tabular}{lrr}
    \toprule
    \textbf{Class} & \textbf{Raw windows} & \textbf{After balancing} \\
    \midrule
    Baseline & 866   & 442 \quad(undersampled) \\
    Sweat    & 160   & 442 \quad($+282$ bootstrapped) \\
    Blood    & 442   & 442 \quad(unchanged) \\
    \midrule
    \textbf{Total} & \textbf{1{,}468} & \textbf{1{,}326} \\
    \bottomrule
\end{tabular}
\end{table}

% ---------------------------------------------------------------- Performance
\section*{Model Performance}

Across model families, the gradient-boosted trees were the strongest. Table~4 compares
overall accuracy on the four held-out sessions. XGBoost with importance-ranked feature
reduction (top 80\% of features) is the best base model.

\begin{table}[H]
\centering
\caption{Overall accuracy on held-out test sessions (sweat, pure blood, and
20\% / 50\% mixtures), before temporal smoothing.}
\begin{tabular}{lcc}
    \toprule
    \textbf{Model} & \textbf{Overall accuracy} & \textbf{Notes} \\
    \midrule
    Random Forest (all features)     & 96.7\% & 200 trees, balanced classes \\
    XGBoost (naive)                  & 97.7\% & all features, defaults \\
    \textbf{XGBoost (20\% feature reduction)} & \textbf{98.6\%} & importance-ranked pruning \\
    1D CNN (regularized)             & 89.4\% & level signal lost in scaling \\
    \bottomrule
\end{tabular}
\end{table}

\subsection*{Best Model: XGBoost with Real-Time Smoothing}

Because each 60\,s window is classified on its own, the raw predictions can flicker between
classes from one window to the next. Since the windows form a time series, a causal
temporal smoother averages the predicted probabilities over the last 5 windows before
deciding, using only current and past data so it runs in real time. This produces stable
predictions suitable for live deployment.

On the held-out sessions the smoothed real-time model reaches \textbf{95.8\% overall}
(Table~5). The small gap from the unsmoothed 98.6\% is onset lag: the smoother takes a few
windows to commit after a sample is introduced. Almost all remaining errors fall in that
onset window; once the response is established, predictions are correct and stable.

\begin{table}[H]
\centering
\caption{Per-session accuracy of the deployed model (XGBoost, 20\% feature reduction,
causal real-time smoothing).}
\begin{tabular}{lcr}
    \toprule
    \textbf{Session} & \textbf{Type} & \textbf{Accuracy} \\
    \midrule
    \texttt{blood\_9}      & Pure blood          & 97.0\% \\
    \texttt{sweat\_7}      & Pure sweat          & 96.6\% \\
    \texttt{mix\_75\_75\_3}  & Mixture, 50\% blood & 95.7\% \\
    \texttt{mix\_120\_30\_4} & Mixture, 20\% blood & 93.0\% \\
    \midrule
    \textbf{Overall} & (215 windows) & \textbf{95.8\%} \\
    \bottomrule
\end{tabular}
\end{table}

The figures below show the model's prediction over each held-out session. Green marks
windows the model classified correctly, red marks errors. The error band sits at the sample
introduction in every session; the long baseline and response regions are predicted
correctly.

\begin{figure}[H]
    \centering
    \begin{minipage}{0.49\textwidth}\centering
        \includegraphics[width=\textwidth]{pred_blood_9.png}
    \end{minipage}\hfill
    \begin{minipage}{0.49\textwidth}\centering
        \includegraphics[width=\textwidth]{pred_sweat_7.png}
    \end{minipage}

    \vspace{0.6em}

    \begin{minipage}{0.49\textwidth}\centering
        \includegraphics[width=\textwidth]{pred_mix_75_75_3.png}
    \end{minipage}\hfill
    \begin{minipage}{0.49\textwidth}\centering
        \includegraphics[width=\textwidth]{pred_mix_120_30_4.png}
    \end{minipage}
    \caption{Real-time predictions per held-out session (green = correct, red = incorrect).
    Errors are confined to the onset transition when the sample is first introduced.}
\end{figure}

\end{document}
